\documentclass[11pt]{amsart}

\usepackage[T1]{fontenc}
\usepackage[utf8]{inputenc}
\usepackage{lmodern}
\usepackage{microtype}
\usepackage{amsmath,amssymb,amsthm}
\usepackage{booktabs}
\usepackage[hidelinks]{hyperref}

\hypersetup{
  pdftitle={
    Verification of the Weighted Szeged Tree-Minimizer Conjecture
    Through Order Eleven
  }
}

\newtheorem{theorem}{Theorem}

\newcommand{\wSz}{\operatorname{wSz}}
\newcommand{\G}{\mathcal{G}}
\newcommand{\T}{\mathcal{T}}
\newcommand{\N}{\mathcal{N}}

\title[The weighted Szeged tree-minimizer conjecture]{
  Verification of the Weighted Szeged Tree-Minimizer Conjecture
  Through Order Eleven
}
%% No author block: author identity is deliberately omitted.
\date{}

\subjclass[2020]{05C35, 05C12, 05C30, 05C85}
\keywords{
  weighted Szeged index, extremal graph, exhaustive enumeration,
  graph6, nauty
}

\begin{document}

\begin{abstract}
Conjecture~2 of Bok, Furtula, Jedli\v{c}kov\'a, and
\v{S}krekovski asserts that, among connected graphs of a fixed order,
a tree attains the minimum weighted Szeged index. For each
$3\le n\le 11$, we evaluate the invariant on one representative of
every isomorphism class of connected simple $n$-vertex graph, totaling
$1{,}018{,}690{,}327$ graphs. The conjecture holds throughout this
range. At $n=3$, the path $P_3$ and the triangle $K_3$ tie at value
$12$; for every $4\le n\le 11$, every minimizer is a tree. We also
determine the exact minimum over non-trees at each order. This is a
finite verification and does not resolve the conjecture for
$n\ge 12$.
\end{abstract}

\maketitle

\section{Result}

All graphs are finite, simple, and connected. For an edge $e=uv$ of
$G$, set
\[
  n_u(e)
  =
  \bigl|
    \{x\in V(G):d_G(x,u)<d_G(x,v)\}
  \bigr|,
\]
and define $n_v(e)$ analogously. The weighted Szeged index is
\[
  \wSz(G)
  =
  \sum_{uv\in E(G)}
  \bigl(\deg_G(u)+\deg_G(v)\bigr)n_u(uv)n_v(uv).
\]
Vertices equidistant from the endpoints of an edge are counted in
neither factor. The index was introduced by Ili\'c and
Milosavljevi\'c~\cite{IlicMilosavljevic2013}. Bok et
al.~\cite{BokEtAl2019} posed the following conjecture.

\medskip

\noindent
\textbf{Weighted Szeged tree-minimizer conjecture
\cite[Conjecture~2]{BokEtAl2019}.}
For $n$-vertex graphs the minimum weighted Szeged index is attained by
a tree.

\medskip

The displayed sentence is the source's wording. Connectivity is
implicit in it: $n_u(e)$ is defined through graph distances, and the
computer tests reported in \cite{BokEtAl2019} range over connected
graphs, so we read the statement over connected graphs and take every
minimum below over that class. The source poses two conjectures. Its
Conjecture~1, that the \emph{maximum} weighted Szeged index over
connected graphs of a given order is attained by the balanced complete
bipartite graph, was proved after submission by
Cambie~\cite{Cambie2021}, as a remark in the published version of
\cite{BokEtAl2019} records. The minimization statement displayed above
is the one at issue here, and it is the one that remains open.

The cases $n=1,2$ are immediate, so the census begins at $n=3$, the
first order admitting a connected non-tree. Let $\G_n$ be the set of
isomorphism classes of connected simple graphs on $n$ vertices, let
$\T_n\subseteq\G_n$ be the trees, and let
$\N_n=\G_n\setminus\T_n$.

\begin{theorem}\label{thm:main}
For every $3\le n\le 11$,
\[
  \min_{G\in\G_n}\wSz(G)
  =
  \min_{T\in\T_n}\wSz(T).
\]
At $n=3$, both $P_3$ and $K_3$ attain the minimum $12$. For every
$4\le n\le 11$, every minimizer is a tree. The exact tree and
non-tree minima are given in Table~\ref{tab:minima}.
\end{theorem}

\begin{table}[ht]
\centering
\small
\caption{Minimum weighted Szeged indices from the exhaustive census.}
\label{tab:minima}
\begin{tabular}{@{}rrrrr@{}}
\toprule
$n$
  & connected classes
  & tree minimum
  & non-tree minimum
  & gap \\
\midrule
3  & $2$                   & $12$  & $12$  & $0$  \\
4  & $6$                   & $34$  & $36$  & $2$  \\
5  & $21$                  & $72$  & $76$  & $4$  \\
6  & $112$                 & $130$ & $132$ & $2$  \\
7  & $853$                 & $204$ & $212$ & $8$  \\
8  & $11{,}117$            & $306$ & $316$ & $10$ \\
9  & $261{,}080$           & $432$ & $444$ & $12$ \\
10 & $11{,}716{,}571$      & $578$ & $600$ & $22$ \\
11 & $1{,}006{,}700{,}565$ & $762$ & $778$ & $16$ \\
\bottomrule
\end{tabular}
\end{table}

The tree minima agree with those reported by Hell,
Hern\'andez-Cruz, and Hosseini~\cite{HellEtAl2024}. The census
additionally determines the exact minimum over non-trees throughout
the completed range.

\section{Exhaustive computation}

We generated one representative of every isomorphism class in
$\G_n$ with \texttt{geng} from the nauty and Traces
package~\cite{McKayPiperno2014}. For each graph, breadth-first search
from every vertex produced the distance matrix. The defining sum was
then evaluated in exact integer arithmetic. Because every generated
graph was connected, it was classified as a tree precisely when it
had $n-1$ edges.

The order-$10$ census was completed as a single stream, as $37$
disjoint edge-count shards, and with residue-class shardings modulo
$64$ and $256$. The order-$11$ census was repeated with
residue-class shardings modulo $240$ and $512$. The complete runs
used two separately written graph6 decoders and two distance kernels,
repeated breadth-first search and Floyd--Warshall. All runs returned
the same graph counts, minimum values, and displayed witnesses.

As external checks, the numbers of connected graphs and trees agree
with OEIS A001349 and A000055, respectively. At order $10$, the $37$
edge-count shard sizes agree term by term with OEIS
A054924~\cite{OEIS}. The evaluator also reproduces the published tree
minima for every $3\le n\le 11$.

A further check compares the objects rather than the values. At the
orders $7\le n\le 11$, where our range meets the minimum-$\wSz$ trees
published in \cite[Table~1]{BokEtAl2019}, the tree witnesses of
Table~\ref{tab:witnesses} are isomorphic to the trees displayed there:
at $n=7,8,9$ they are the spiders $S(2,2,2)$, $S(2,2,3)$ and
$S(2,2,2,2)$, and at $n=10,11$ each consists of two adjacent vertices
of degree $3$ carrying legs of length $2$, with a single leg of
length $3$ at $n=11$. These structures are readable directly from the
graph6 strings.

\begin{proof}[Proof of Theorem~\ref{thm:main}]
The census evaluates one representative of every isomorphism class
in $\G_n$, and $\wSz$ is invariant under isomorphism. Hence the two
restricted minima are the entries of Table~\ref{tab:minima}. Direct
evaluation gives
\[
  \wSz(P_3)=\wSz(K_3)=12.
\]
For $4\le n\le 11$, the minimum over $\N_n$ is strictly larger than
the minimum over $\T_n$, so every global minimizer is a tree.
\end{proof}

\section{Witnesses and scope}\label{sec:witnesses}

Table~\ref{tab:witnesses} gives a graph6 witness for the tree minimum
at every completed order, and for the non-tree minimum at the two
largest orders. These records certify the displayed upper bounds
directly from the definition; optimality follows from the exhaustive
census. At $n=6$ the tree minimum is attained by two isomorphism
classes and one of them is listed. Every string in the table was
re-decoded, re-encoded byte for byte, and re-evaluated by both
distance kernels.

\begin{table}[ht]
\centering
\small
\caption{Graph6 witnesses attaining the minima of
Table~\ref{tab:minima}.}
\label{tab:witnesses}
\begin{tabular}{@{}cllr@{}}
\toprule
$n$ & class & graph6 & $\wSz$ \\
\midrule
3  & tree     & \texttt{Bo}           & $12$  \\
4  & tree     & \texttt{Cq}           & $34$  \\
5  & tree     & \texttt{DkG}          & $72$  \\
6  & tree     & \texttt{Eh\_O}        & $130$ \\
7  & tree     & \texttt{FCOf?}        & $204$ \\
8  & tree     & \texttt{GiE@?O}       & $306$ \\
9  & tree     & \texttt{HiQ@?\_G}     & $432$ \\
10 & tree     & \texttt{I?AA@\_gw?}   & $578$ \\
10 & non-tree & \texttt{I?AA@\_gwO}   & $600$ \\
11 & tree     & \texttt{J?AA@AOEB\_?} & $762$ \\
11 & non-tree & \texttt{J?AACGoIF??}  & $778$ \\
\bottomrule
\end{tabular}
\end{table}

Theorem~\ref{thm:main} is a finite computer-assisted verification. It
does not settle the conjecture for $n\ge 12$. Bok et al.\ publish the
minimum-$\wSz$ trees for the orders $7$ to $25$
\cite[Table~1]{BokEtAl2019}; the computer tests supporting
Conjecture~2, namely those ranging over all connected graphs, are
reported without stating the orders covered, and that census is not
published. The statement was still recorded as open after those tests:
Hell, Hern\'andez-Cruz, and Hosseini restate it as their
Conjecture~1.1 and describe characterizing the graphs of fixed order
with minimum weighted Szeged index as one of the main open problems in
the area~\cite{HellEtAl2024}.

\medskip

\noindent
\textbf{Exact verification.}
Every number printed above that can be recomputed from the data of
this note has been recomputed independently, by a second program
written from the statements here rather than from the census code,
using a standard library only and exact integer arithmetic. It
re-derives all eleven witnesses of
Table~\ref{tab:witnesses} and their indices by two distance kernels
and a third edge-split kernel for trees, together with their
byte-exact graph6 re-encodings, and confirms by isomorphism that the
tree witnesses at $7\le n\le 11$ have the structures described above;
the tie
$\wSz(P_3)=\wSz(K_3)=12$ at $n=3$; the whole tree-minimum column for
$3\le n\le 11$, by exhaustive free-tree enumeration; both minima for
$3\le n\le 9$, by exhaustive generation for $n\le 8$ and a covering
enumeration at $n=9$; the connected-class column and its total
$1{,}018{,}690{,}327$, from a cycle-index count and the inverse Euler
transform independently of OEIS; and the gap column. It reports $83$
checks and all $83$ pass. Its non-tree search at $n=10$ and $n=11$ was
exhaustive only over the connected graphs with at most $13$ edges,
where it returns $600$ and $778$, and the bound
$\wSz(G)\ge\sum_{v}\deg(v)^2$ closes $m\ge 39$ at $n=10$ and
$m\ge 47$ at $n=11$; the intermediate edge counts require
\texttt{geng}, and the shard replications were not repeated, so the
non-tree minima at those two orders, the last two gaps, and the
conclusion that every minimizer is a tree there rest on the census
reported above. The program and its transcript accompany this note.

\section*{Reproducibility}

The census uses nothing beyond \texttt{geng} of~\cite{McKayPiperno2014}
and an evaluator of the defining sum, so the exact invocations and the
evaluator are printed here in full. One representative of
every isomorphism class in $\G_n$ is generated as a single stream, as
one shard per edge count, and as residue shards, by

\begin{verbatim}
geng -c -q  n       # every connected class on n vertices
geng -c -q  n  m:m  # one shard per edge count m,
                    #   m = n-1,...,n(n-1)/2
                    #   (37 shards at n = 10)
geng -c -q  n  r/s  # one residue shard r = 0,...,s-1,
                    #   s = 64, 256  (n = 10)
                    #   s = 240, 512 (n = 11)
\end{verbatim}

\noindent
and each graph6 line of the output is evaluated by

\begin{verbatim}
read n and the adjacency of the graph from the graph6 line
for s in 0..n-1:  d[s][*] <- BFS(s)   # or Floyd-Warshall
w <- 0
for each edge uv:
    nu <- #{x : d[x][u] < d[x][v]}
    nv <- #{x : d[x][v] < d[x][u]}
    w  <- w + (deg(u) + deg(v)) * nu * nv
if (number of edges) = n-1:
    update (tree minimum, witness) with (w, line)
else:
    update (non-tree minimum, witness) with (w, line)
\end{verbatim}

\noindent
All arithmetic is exact integer arithmetic; there is no randomization,
no floating point, and no tolerance parameter. Every entry of
Tables~\ref{tab:minima} and~\ref{tab:witnesses} is an output of this
loop, and each witness can be rechecked against the definition from
its graph6 string alone. No data file accompanies this note, and none
is needed to recheck it: the generated shard files are too
large to deposit and are regenerated by the commands above rather than
reproduced, and the independent program of
Section~\ref{sec:witnesses}, which accompanies this note together with
its transcript, takes no input beyond the tables printed here.

\begin{thebibliography}{9}

\bibitem{IlicMilosavljevic2013}
A.~Ili\'c and N.~Milosavljevi\'c,
\emph{The weighted vertex PI index},
Math. Comput. Modell. \textbf{57} (2013), 623--631,
\href{https://doi.org/10.1016/j.mcm.2012.08.001}
{doi:10.1016/j.mcm.2012.08.001}.

\bibitem{BokEtAl2019}
J.~Bok, B.~Furtula, N.~Jedli\v{c}kov\'a, and R.~\v{S}krekovski,
\emph{On extremal graphs of weighted Szeged index},
MATCH Commun. Math. Comput. Chem. \textbf{82} (2019), 93--109,
\href{https://arxiv.org/abs/1901.04764}{arXiv:1901.04764}.

\bibitem{Cambie2021}
S.~Cambie,
\emph{Five results on maximizing topological indices in graphs},
Discrete Math. Theor. Comput. Sci. \textbf{23} (2021), no.~3,
Paper No.~dmtcs:6896,
\href{https://doi.org/10.46298/dmtcs.6896}{doi:10.46298/dmtcs.6896}.

\bibitem{HellEtAl2024}
P.~Hell, C.~Hern\'andez-Cruz, and S.~A.~Hosseini,
\emph{Trees with minimum weighted Szeged index},
Art Discrete Appl. Math. \textbf{7} (2024), no.~2,
Paper No.~P2.09,
\href{https://doi.org/10.26493/2590-9770.1702.63d}
{doi:10.26493/2590-9770.1702.63d}.

\bibitem{McKayPiperno2014}
B.~D.~McKay and A.~Piperno,
\emph{Practical graph isomorphism, II},
J. Symbolic Comput. \textbf{60} (2014), 94--112,
\href{https://doi.org/10.1016/j.jsc.2013.09.003}
{doi:10.1016/j.jsc.2013.09.003}.

\bibitem{OEIS}
OEIS Foundation Inc.,
\emph{The On-Line Encyclopedia of Integer Sequences},
entries A001349, A000055, and A054924,
\url{https://oeis.org}.

\end{thebibliography}

\end{document}